\chapter{Conclusion and Future Work}\label{ch:future_work}

\section{Conclusion}

This thesis investigated how a cell-free ISAC transmitter can minimize total
transmit power while providing robust communication QoS and target-specific
sensing service with a physically implementable AP--target topology. The
resulting design keeps communication beamforming globally cooperative, while
binary association variables authorize only the dedicated sensing covariance
of each target. This separation aligns the sensing topology, the
target-specific FIM/PCRB requirement, and the dedicated sensing power directly
controlled by the optimization.

The formulated mixed-integer problem jointly accounts for worst-case
communication SINR, sensing SINR, matrix PCRB, dedicated-sensing
authorization, and total AP-power limits. Covariance lifting, the
S-Procedure, and a Schur-complement representation yield solver-ready robust
and tracking constraints. A dual DC-penalty SCA procedure then reduces both
rank-one and binary deficiencies, and direct validation distinguishes a
recovered physical design from a relaxed optimization point.

The numerical results support three bounded conclusions. First,
topology-aware recovery improves physical feasibility relative to nearest-AP
and random topology rules in the tested ordinary and pressure geometries.
Second, robust design incurs a modest transmit-power increase while markedly
reducing observed outage within the specified CSI-uncertainty model. Third,
the dual-DC workflow provides recovery-quality evidence through small rank and
binary residuals before physical validation. The larger-dimensional experiment
also demonstrates usable recovery behaviour at the tested scale, at a higher
computational cost than fixed-topology baselines. These are empirical,
configuration-specific findings; they do not establish global optimality,
universal relaxation tightness, or asymptotic scalability.

\section{Future Work}

The immediate priority is broader, reproducible validation. Future experiments
should report feasibility, recovered physical transmit power, worst-case
communication SINR, sensing SINR, PCRB, per-AP power, and
association-cardinality residuals for every sweep, while preserving the
distinction between relaxed and recovered solutions. Larger independent seed
sets, confidence intervals, and systematic PCRB, cluster-size, and CSI-error
sweeps are needed to quantify the robustness--sensing--power trade-off more
reliably.

The model can also be extended beyond its single-interval assumptions through
dynamic target-state prediction across scheduling intervals, uncertainty in
sensing channels and target parameters, fronthaul or processing costs for
sensing clusters, and distributed implementations for larger cell-free
networks. Each extension should retain direct validation of the recovered
physical variables so that additional modelling flexibility does not obscure
the communication, sensing, tracking, and AP-power guarantees.
